\chapter{\bilingual{Implementación}{Implementation}}

\section{\bilingual{Estructura de la aplicación}{Application Structure}}

\bilingual{%
  Describe aquí la organización del código fuente (módulos, capas,
  paquetes\ldots).
}{%
  Describe here how the source code is organized (modules, layers,
  packages\ldots).
}

% See ch02's note: \bilingual must not be used directly inside a caption
% (here, the lstlisting "caption=" key), so resolve it to a plain macro
% first.
\ifthesisenglish\def\ExampleListingCaption{Example code snippet}\else\def\ExampleListingCaption{Ejemplo de fragmento de código}\fi
\begin{lstlisting}[language=Python, caption={\ExampleListingCaption}, label={lst:example}]
def average(values):
    """Return the arithmetic mean of a list of values."""
    return sum(values) / len(values)
\end{lstlisting}

\bilingual{%
  Como se observa en el Código~\ref{lst:example}, \ldots
}{%
  As shown in Listing~\ref{lst:example}, \ldots
}

\bilingual{%
  El mismo ejemplo, ahora en Java, importado directamente desde su propio
  fichero fuente en \texttt{snippets/} (en vez de escribirlo dentro del
  \texttt{.tex}) con \texttt{\textbackslash{}lstinputlisting}:
}{%
  The same example, now in Java, imported directly from its own source
  file in \texttt{snippets/} (instead of being typed inside the
  \texttt{.tex} file) with \texttt{\textbackslash{}lstinputlisting}:
}

\ifthesisenglish\def\JavaListingCaption{Example code snippet (Java)}\else\def\JavaListingCaption{Ejemplo de fragmento de código (Java)}\fi
\lstinputlisting[language=Java, caption={\JavaListingCaption}, label={lst:example-java}]{snippets/Average.java}

\section{\bilingual{Implementación de las pruebas}{Test Implementation}}

\bilingual{%
  Describe aquí cómo se han implementado las pruebas definidas en el
  capítulo de diseño.
}{%
  Describe here how the tests defined in the design chapter were
  implemented.
}
